% Auto-generated from 44-Convergence.md
% POV: Aisen | Landscape: Coast | Location: The Pilgrim's Path

By month sixteen of the overall conflict cycle, provincial resistance networks, Academy network operations, and military officers supporting resistance interests had converged into something resembling a coordinated structure. It wasn't unified—there were still separate cells, competing strategic philosophies, different regional priorities. But there was coordination.

Aisen, working through Kael's coordinated network operations, was attempting to build framework where the different resistance structures could function in sufficient coordination to challenge Covenant institutional authority effectively without requiring complete ideological alignment or unified command structure.

The key insight was understanding that the Covenant's institutional control required consensus from multiple institutional sectors—education was one component, governance was another, military maintenance was a third, economic activity was a fourth. If resistance networks could position themselves to subtly challenge Covenant control within multiple sectors simultaneously, institutional authority would face pressure it couldn't address through simple security crackdown.

"We're not looking for rapid institutional change," Aisen said to coordinated network leadership during a rare full assembly meeting. "We're looking for gradual institutional erosion. We want Covenant administrative authority to become sufficiently expensive to maintain that alternatives start appearing possible."

"How long does that timescale extend?" someone asked.

"Years, possibly decades," Aisen said. "But it's strategy designed to create conditions where provincial autonomy could eventually reemerge through civil society mechanisms rather than through military resistance."

"How do we maintain operational coherence across that timeframe?" Kael asked.

"By keeping resistance networks focused on specific institutional sector objectives," Aisen said. "Military liaison networks maintain communication and tactical coordination. Educational networks focus on curriculum influence. Governance networks focus on administrative procedure influence. Economic networks focus on resource distribution alternatives. As long as each sector maintains its own institutional operation, the network doesn't need complete unified command structure—just sufficient coordination to avoid stepping on each other's work."

The approach was unconventional but strategically sound. The Covenant's institutional control was maintained through distributed structures that didn't require perfect coordination to function. Resistance networks could similarly function through distributed structures that maintained coordination at strategic level while preserving operational independence at tactical level.
